\section{Capping primitive annuli without another surface search}
\label{sec:cocycle-annulus}

The optimal-face experiment exposes a simpler improvement than searching
for a tied disc: its seven-crossing source already has a certified annulus
in the primitive relative class. That annulus itself proves unknottedness.
The native stage now accepts this sufficient witness immediately. The
stored surface remains an annulus with zero disc components; a new
certificate schema records the separate boundary-capping implication.

\subsection{Why the primitive annulus gives a compressing disc}

Let $M$ be the authenticated compact exterior of a knot in $S^3$, with
boundary torus $T$. As used in Section~\ref{sec:cocycle-connectivity},
\[
 H_2(M;\mathbb Z)=0,\qquad H_2(M,T;\mathbb Z)\cong\mathbb Z,
 \qquad H_1(M;\mathbb Z)\cong\mathbb Z.
\]
The inclusion $i_*:H_1(T;\mathbb Z)\to H_1(M;\mathbb Z)$ is onto.
For example, Mayer--Vietoris for $S^3=M\cup N(K)$ identifies
$H_1(T;\mathbb Z)$ with $H_1(M;\mathbb Z)\oplus H_1(N(K);\mathbb Z)$.
The exact sequence of the pair therefore makes
\[
 \partial:H_2(M,T;\mathbb Z)\longrightarrow H_1(T;\mathbb Z)
\]
injective, with image the primitive rank-one subgroup $\ker i_*$.
Consequently a primitive relative class has nonzero primitive boundary
class on the torus. This conclusion uses the source knot exterior; it is
not asserted for an arbitrary torus-boundary manifold.

Suppose $S$ is a connected orientable proper surface representing that
primitive class, with $\chi(S)=0$. Classification of compact orientable
surfaces leaves a torus or an annulus. A closed torus has zero boundary
class, which contradicts the preceding paragraph. Thus $S$ is an annulus.

Its two boundary circles are disjoint simple curves on $T$. If both were
inessential, their homology sum would be zero. If both were essential,
they would be parallel on the torus; their oriented classes would sum to
zero or twice a primitive class. Neither possibility is nonzero primitive.
Exactly one boundary circle is therefore essential and the other bounds
a disc $D\subset T$. The essential circle lies outside that disc: an
essential curve cannot lie inside a disc and cannot cross its boundary.

The interior of the proper annulus is disjoint from $T$, so $S\cup D$
is an embedded disc with boundary the remaining essential circle. Using
a boundary collar, push its boundary cap inward and round the corner
along the capped circle. The result is a properly embedded compressing
disc for $T$. The knot-exterior disc criterion then proves that the source
knot is the unknot. No normal vector for this capped disc is needed to
verify the implication.

The two essential hypotheses cannot be replaced by Euler characteristic
alone. The audit below includes a connected null-class annulus in a trefoil
exterior and a primitive disconnected surface there with total Euler
characteristic zero. Both are rejected. Coorientation is also essential:
the connectedness inspector proves orientability rather than inferring it
from Euler characteristic.

\subsection{A separate witness tag and unchanged surface inspection}

The new schema is \code{diagram-cocycle-annulus-v1}. It contains the same
six fields as the previous cocycle disc schema: source PD, canonical finite
triangulation, integral heights, normal coordinates, optional span witness,
and schema tag. The new independent wrapper translates only the tag for
the existing surface inspector and requires Euler characteristic exactly
zero. That inspector still checks the source geometry, all matching and
coherence conditions, the primitive class after its own tree gauge, and
either a zero spanning tree or the minimum-span certificate. Its code and
the span optimality verifier are unchanged.

The public seed verifier dispatches the new schema to this wrapper.
An annulus relabeled as the old disc schema remains invalid as a positive
disc proof; a disc relabeled as an annulus also fails. Legacy cocycle-disc
and component-count certificates retain their previous acceptance rules.
In particular, the new witness does not redefine a zero disc count as a
positive normal-disc count.

Discovery already performs independent surface inspection at the raw and
optimized stages. A verified Euler-zero result now stops that search and
emits the annulus tag. Tree and face candidates with Euler zero must also
pass full independent source inspection before being accepted; the cheap
negative-candidate branch cannot emit such a positive. The final stage
record has \code{unknot\_witness=annulus-cap} and continues to report
\code{compressing\_discs=0} for its actual stored surface.

The new implication adds only constant-size metadata and a constant number
of checks after surface inspection. The standalone wrapper repeats the
same independent inspection and adds the Euler-zero test. It introduces
no new search, no boundary-component enumeration, and no expansion of
binary normal coordinates. All previously accounted-for arithmetic,
container and bit costs of source inspection remain applicable. The
capping argument is an existence certificate for a disc; it does not yet
provide a normal-coordinate disc extraction or a hierarchy-cut operation.

Annulus recognition is enabled within the optional native stage.
The seed routine accepts \code{annulus=False} to recover the old criterion.
The recognizer exposes \code{normal\_seed\_annulus=False}; the CLI exposes
\code{--normal-seed-no-annulus}. The shared
work allowance and cancellation behavior are retained. The separate face
search remains disabled by default, and the ordinary recognition portfolio
still leaves the whole native stage optional.

\subsection{Coverage and independent boundary controls}

On 84 source inputs with the existing two-million-guard allowance, the
four-tree policy now has 24 native positives, 59 completed restricted
misses and one cap, compared with 23 positives before annulus recognition.
Requesting 24 tree trials now gives 27 positives, 56 misses and one cap,
compared with 24 positives before. The capped source remains Gordian.
The optional four-face-root policy retains its 24 positives; its newly
found disc is now preempted by the earlier optimized annulus.

Three earlier positives, \code{random-29}, \code{random-79} and
\code{random-99}, now stop at their raw annuli. Their shared guard work
falls from 11,227, 10,192 and 7,785 to 6,561, 6,480 and 4,984 respectively.
The seven-crossing \code{random-142} stops on its original 122-piece
optimized annulus at 93,012 guards. Without annulus recognition and without
face search it previously completed a miss at 93,011 guards: the extra
unit is the final success guard, not another discovery pass. With four
face roots enabled, its old disc proof required 117,848 guards.

With the optional 24-tree allowance, \code{random-81} now succeeds at
trial 5 rather than trial 22, reducing guards from 72,527 to 52,641.
Two additional four-crossing unknots, \code{random-1} and
\code{random-127}, succeed at tree trials 10 and 7 respectively. Their
annuli have 86 and 111 pieces, and their guard totals are 44,244 and 48,849.
These are additional positives for that explicit extended allowance;
they remain misses under the default four tree trials.

Seven distinct annulus certificates are retained. Every prior positive
certificate still passes the public verifier. Under each of the three
compared policies, twenty positive records retain their exact certificate
hash; the intentional changes use the new annulus tag and may stop earlier.
Cases that do not return an annulus preserve their complete recorded native
result, including their work count. On nineteen ordinary corpus inputs,
both settings return the expected verdicts without reaching the native
stage, so no ordinary-portfolio gain follows from this experiment.

The final audit independently checks all 787 original or root-extremal
minimum-span surfaces from the preceding 74-source small-input experiment.
Sixteen have Euler characteristic zero, and precisely those pass the new
annulus wrapper. Regina agrees on connectedness, orientability, Euler
characteristic and boundary count. The seven actual produced annulus
witnesses receive separate Regina comparisons, including raw and later-tree
surfaces outside that minimum-span experiment. Two trefoil adversaries
bring the total to 796 Regina comparisons.

For the sixteen accepted minimum-span annuli and the seven produced
witnesses, a separate small-fixture control expands boundary arc pairings,
finds each boundary circle and computes its mod-two intersection vector
against an independently validated torus homology basis. All 23 comparisons
give two circles, exactly one with zero vector. At most 86 boundary points
are expanded. This is audit code only; neither recognition nor certificate
verification calls it.

The first trefoil adversary is a genuine connected orientable annulus with
20 normal pieces. Set the local height at a corner to one exactly when its
canonical global vertex is 3 or 7, and to zero otherwise. This is a pure
vertex coboundary. Its two boundary circles both have zero mod-two homology,
and the primitive-class checker rejects it. The second adversary changes
the primitive trefoil seed by subtracting one at global vertex 3. Its total
Euler characteristic is zero, but Regina and the native component census
find two components and no compressing disc. The zero-tree connectedness
check rejects this forged annulus witness.

\subsection{Tests and complete-call measurements}

All 34 focused tests pass in 4.809 seconds, and all 1,204 maintained tests
pass in 226.376 seconds. Tests cover producer-free annulus replay, raw and
optimized early termination, exact budget boundaries, cancellation, strict
switch types, schema swaps, malformed fields, wrong sources, the trefoil
adversaries, zero and doubled classes, and 20,000-bit gauge transport through
hexadecimal JSON. Historical tree and face tests explicitly disable the new
criterion when testing their original discovery paths.

The final audit takes 159.502 seconds and pins 488 source files. Its
elapsed time is a correctness-run total, not a paired performance result.
A preliminary successful audit is retained as a log; the final run follows
an update that makes the historical tree, blocking-flow and face research
drivers explicitly retain their original disc-only comparison criterion.
The cocycle-surface inspector and span verifier remain byte-identical to
baseline \code{b2d13b5f4}.

Measurements run serially without overlapping tests, audits, builds or
commits. They start from the input diagram and include the complete
recognizer, exact fallback when needed, independent replay of any returned
native proof and certificate hashing. Earlier shortcuts are disabled to
reach the native stage. Both arms have no local seed cap and retain the
twelve-second global deadline. The on/off comparison changes only the
annulus criterion. Tree trials are four except for the three explicitly
extended-tree cases; face roots are zero except for the labeled prior
face-search positive, which requests four.

The initial five-round run has visible A/A variation, including a 1.089
raw-disc on-arm ratio and a 0.876 ratio for the three-crossing raw annulus.
It is retained, and a nine-round repeat supplies the final table. Each run
has one warm-up round per case, and each round randomizes off, off A/A,
on and on A/A endpoints. All 616 measured calls and 88 warm-ups across the
two runs finish. Their per-implementation verdicts, witness hashes and
native work counts are stable across both runs. The final run has 396
measured calls and 44 warm-ups and takes 37.153 seconds.

\begin{center}
\small\begin{tabular}{lrrrrr}
\toprule
Input & Off (ms) & On (ms) & Off/on & Off A/A & On A/A \\
\midrule
Raw disc control & 64.707 & 65.940 & 0.992 & 0.958 & 0.990 \\
Optimized disc control & 91.996 & 90.158 & 1.012 & 1.008 & 0.995 \\
Raw annulus, 4 crossings A & 50.989 & 31.007 & 1.637 & 1.001 & 0.995 \\
Raw annulus, 4 crossings B & 49.219 & 31.143 & 1.580 & 0.998 & 0.999 \\
Raw annulus, 3 crossings & 36.943 & 23.334 & 1.571 & 0.997 & 1.008 \\
New optimized annulus & 118.295 & 137.411 & 0.865 & 1.005 & 0.991 \\
Prior face-search positive & 206.528 & 137.443 & 1.483 & 0.998 & 0.988 \\
Earlier late-tree positive & 136.703 & 100.128 & 1.349 & 1.031 & 0.970 \\
New tree annulus A & 89.212 & 85.510 & 1.033 & 1.021 & 1.004 \\
New tree annulus B & 92.456 & 85.370 & 1.075 & 0.993 & 0.998 \\
Trefoil miss & 39.753 & 40.171 & 0.992 & 0.998 & 1.001 \\
\bottomrule
\end{tabular}

\end{center}

The three early-annulus cases improve $1.57$--$1.64\times$ by paired
ratios. The previous face-search positive improves $1.48\times$, and the
late-tree positive improves $1.35\times$. The two new extended-tree
positives have ratios 1.033 and 1.075; the first is modest relative to its
A/A variation, so the main conclusion there is added native coverage.
The unchanged raw-disc, optimized-disc and trefoil controls are near parity.
Table times are primary-arm medians and ratios are medians of paired
ratios; these need not be quotients of the displayed median times.

The newly recognized seven-crossing source without face search is an
adverse complete-call control. Its old exact fallback finishes in a median
118 milliseconds; returning and independently replaying the new native
proof takes 137 milliseconds, about 16 percent slower by the paired ratio.
There is no claim of a uniform recognition speedup. The criterion adds no
surface discovery loop and shortens several existing native searches, but
native proof production and replay can cost more than the old fallback on
a small diagram.

The reproduction driver is \path{../fast/normal_orbit_research/annulus.py}.
Use its \code{audit} or \code{recognize --rounds 9} subcommand;
the audit uses optional Regina.
The incoming archives and checked \code{openai/math} head
\code{fd4aeeb2e} are unchanged. This is a stronger polynomial sufficient
certificate test within an incomplete candidate search, not a general
subexponential or quasipolynomial recognition algorithm.

